\documentclass[11pt, a4paper]{article}
\usepackage[a4paper, margin=2cm]{geometry}
\usepackage{fontspec}
\usepackage[english, russian, bidi=basic, provide=*]{babel}
\babelprovide[import, onchar=ids fonts]{russian}
\babelfont{rm}{Noto Sans}

\usepackage{amsmath}
\usepackage{amssymb}
\usepackage{booktabs}
\usepackage{multicol}
\usepackage{xcolor}
\usepackage[colorlinks=true, linkcolor=blue, urlcolor=blue]{hyperref}

% Определение цветов для оформления
\definecolor{sectioncolor}{RGB}{0, 51, 102}
\definecolor{subsectioncolor}{RGB}{34, 112, 147}

\title{\textbf{\Large Справочник по физике: Основные формулы и определения}}
\author{Сводный курс формул}
\date{\today}

\begin{document}

\maketitle

\tableofcontents
\newpage

% ==================== 1. КИНЕМАТИКА ====================
\section{Кинематика}

\subsection{Равномерное прямолинейное движение}
\begin{itemize}
    \item \textbf{Путь:} $L = v \cdot t$ \\
    \textit{где} $L$ --- путь, $v$ --- скорость, $t$ --- время.
    
    \item \textbf{Уравнение движения (координата):} $x = x_0 + v_x t$ \\
    \textit{где} $x_0$ --- начальная координата, $v_x$ --- проекция скорости на ось $X$.
    
    \item \textbf{Средняя путевая скорость:} $v_{cp} = \dfrac{L_{\text{полн}}}{t_{\text{полн}}} = \dfrac{L_1 + L_2 + \dots}{t_1 + t_2 + \dots}$
\end{itemize}

\subsection{Равноускоренное прямолинейное движение}
\begin{itemize}
    \item \textbf{Ускорение:} $a = \dfrac{\Delta v}{\Delta t} = \dfrac{v - v_0}{t}$ \\
    \textit{где} $a$ --- ускорение, $v_0$ --- начальная скорость, $v$ --- конечная скорость.
    
    \item \textbf{Скорость от времени:} $v = v_0 + at$
    
    \item \textbf{Перемещение:} $S = v_0 t + \dfrac{at^2}{2} = \dfrac{v^2 - v_0^2}{2a} = \dfrac{v + v_0}{2} t$
    
    \item \textbf{Уравнение координаты:} $x = x_0 + v_{0x} t + \dfrac{a_x t^2}{2}$
    
    \item \textbf{Тормозной путь:} $S_{\text{торм}} = \dfrac{v_0^2}{2a}$
\end{itemize}

\subsection{Свободное падение и движение тела, брошенного вертикально вверх}
\begin{itemize}
    \item \textbf{Падение без начальной скорости с высоты $h$:}
    \[ v = \sqrt{2gh}, \quad t_{\text{пад}} = \sqrt{\dfrac{2h}{g}} \]
    \textit{где} $g \approx 9,8 \text{ м/с}^2$ --- ускорение свободного падения.
    
    \item \textbf{Подъем вертикально вверх с начальной скоростью $v_0$:}
    \[ h_{\max} = \dfrac{v_0^2}{2g}, \quad t_{\text{под}} = \dfrac{v_0}{g}, \quad t_{\text{полн}} = \dfrac{2v_0}{g} \]
\end{itemize}

\subsection{Движение тела, брошенного под углом к горизонту}
\begin{itemize}
    \item \textbf{Время полета и дальность:}
    \[ t_{\text{полн}} = \dfrac{2v_0 \sin\alpha}{g}, \quad L = \dfrac{v_0^2 \sin(2\alpha)}{g}, \quad h_{\max} = \dfrac{(v_0 \sin\alpha)^2}{2g} \]
\end{itemize}

\subsection{Равномерное движение по окружности}
\begin{itemize}
    \item \textbf{Период и частота:} $T = \dfrac{t}{N} = \dfrac{1}{\nu}, \quad \nu = \dfrac{N}{t} = \dfrac{1}{T}$
    
    \item \textbf{Угловая и линейная скорость:} $\omega = \dfrac{\Delta\varphi}{\Delta t} = 2\pi\nu = \dfrac{2\pi}{T}, \quad v = \omega R = \dfrac{2\pi R}{T}$
    
    \item \textbf{Центростремительное ускорение:} $a_n = \dfrac{v^2}{R} = \omega^2 R = v\omega$
\end{itemize}

% ==================== 2. ДИНАМИКА ====================
\section{Динамика}

\subsection{Законы Ньютона}
\begin{itemize}
    \item \textbf{Второй закон Ньютона:} $\vec{F} = m\vec{a}$ \quad или в проекциях: $\sum F_x = m a_x$ \\
    \textit{где} $\vec{F}$ --- равнодействующая всех сил, $m$ --- масса тела.
    
    \item \textbf{Третий закон Ньютона:} $\vec{F}_{12} = -\vec{F}_{21}$
\end{itemize}

\subsection{Силы в механике}
\begin{itemize}
    \item \textbf{Сила упругости (Закон Гука):} $F_{\text{упр}} = k x$ \\
    \textit{где} $k$ --- жесткость, $x$ --- абсолютная деформация.
    \begin{itemize}
        \item Последовательное соединение пружин: $\dfrac{1}{k_{\text{эф}}} = \dfrac{1}{k_1} + \dfrac{1}{k_2}$
        \item Параллельное соединение пружин: $k_{\text{эф}} = k_1 + k_2$
    \end{itemize}

    \item \textbf{Сила трения скольжения:} $F_{\text{тр}} = \mu N$ \\
    \textit{где} $\mu$ --- коэффициент трения, $N$ --- сила реакции опоры.

    \item \textbf{Закон всемирного тяготения:} $F = G \dfrac{M m}{r^2}$ \\
    \textit{где} $G = 6,67 \cdot 10^{-11} \text{ Н}\cdot\text{м}^2/\text{кг}^2$ --- гравитационная постоянная.

    \item \textbf{Ускорение свободного падения на высоте $h$:} $g_h = \dfrac{G M}{(R_n + h)^2}$
    
    \item \textbf{Первая космическая скорость:} $v_1 = \sqrt{g R_n} = \sqrt{\dfrac{G M}{R_n}}$
\end{itemize}

% ==================== 3. СТАТИКА И ГИДРОСТАТИКА ====================
\section{Статика и Гидростатика}

\subsection{Статика}
\begin{itemize}
    \item \textbf{Момент силы:} $M = F \cdot d$ \quad (\textit{условие равновесия}: $\sum M = 0$).
    \item \textbf{Центр масс системы точек:} $x_{cm} = \dfrac{\sum m_i x_i}{\sum m_i}$
\end{itemize}

\subsection{Гидростатика}
\begin{itemize}
    \item \textbf{Давление:} $p = \dfrac{F}{S}$
    \item \textbf{Гидростатическое давление:} $p = \rho g h$
    \item \textbf{Полное давление на глубине $h$:} $p = p_0 + \rho g h$ \quad ($p_0$ --- атмосферное давление).
    \item \textbf{Закон Архимеда (выталкивающая сила):} $F_A = \rho_{\text{жидк}} g V_{\text{погр}}$
\end{itemize}

% ==================== 4. ИМПУЛЬС И ЭНЕРГИЯ ====================
\section{Импульс и Работа. Законы сохранения}

\subsection{Импульс}
\begin{itemize}
    \item \textbf{Импульс тела:} $\vec{p} = m \vec{v}$
    \item \textbf{Второй закон Ньютона в импульсной форме:} $\vec{F} \Delta t = \Delta \vec{p}$
    \item \textbf{Закон сохранения импульса:} $\sum \vec{p}_{\text{нач}} = \sum \vec{p}_{\text{кон}}$ (при $\vec{F}_{\text{внеш}} = 0$).
\end{itemize}

\subsection{Работа и Энергия}
\begin{itemize}
    \item \textbf{Механическая работа:} $A = F S \cos\alpha$
    \item \textbf{Мощность:} $P = \dfrac{A}{t} = F v \cos\alpha$
    \item \textbf{Кинетическая энергия:} $E_k = \dfrac{m v^2}{2} = \dfrac{p^2}{2m}$
    \item \textbf{Потенциальная энергия:} $E_p = m g h$ (в поле тяжести), \quad $E_p = \dfrac{k x^2}{2}$ (упругая деформация).
    \item \textbf{Теорема об изменении механической энергии:} $A_{\text{внеш}} = E_{\text{кон}} - E_{\text{нач}}$
\end{itemize}

% ==================== 5. МОЛЕКУЛЯРНАЯ ФИЗИКА И ТЕРМОДИНАМИКА ====================
\section{Молекулярная физика и Термодинамика}

\subsection{Основы МКТ}
\begin{itemize}
    \item \textbf{Количество вещества:} $\nu = \dfrac{m}{M} = \dfrac{N}{N_A}$
    \item \textbf{Основное уравнение МКТ:} $p = \dfrac{1}{3} n m_0 v_{\text{кв}}^2 = n k T$
    \item \textbf{Средняя кинетическая энергия поступательного движения молекулы:} $\langle \epsilon \rangle = \dfrac{3}{2} k T$
    \item \textbf{Средняя квадратичная скорость:} $v_{\text{кв}} = \sqrt{\dfrac{3 k T}{m_0}} = \sqrt{\dfrac{3 R T}{M}}$
    \item \textbf{Уравнение Менделеева--Клапейрона:} $p V = \nu R T = \dfrac{m}{M} R T$
\end{itemize}

\subsection{Изопроцессы}
\begin{itemize}
    \item \textbf{Изотермический ($T = \text{const}$):} $p V = \text{const}$ (Закон Бойля--Мариотта)
    \item \textbf{Изобарный ($p = \text{const}$):} $\dfrac{V}{T} = \text{const}$ (Закон Гей-Люссака)
    \item \textbf{Изохорный ($V = \text{const}$):} $\dfrac{p}{T} = \text{const}$ (Закон Шарля)
\end{itemize}

\subsection{Термодинамика}
\begin{itemize}
    \item \textbf{Внутренняя энергия одноатомного идеального газа:} $U = \dfrac{3}{2} \nu R T = \dfrac{3}{2} p V$
    \item \textbf{Работа газа при изобарном процессе:} $A = p \Delta V$
    \item \textbf{Первое начало термодинамики:} $Q = \Delta U + A$
    \item \textbf{КПД тепловой машины:} $\eta = \dfrac{A_{\text{полезн}}}{Q_1} = \dfrac{Q_1 - |Q_2|}{Q_1}$
    \item \textbf{Максимальный КПД (цикл Карно):} $\eta_{\max} = \dfrac{T_1 - T_2}{T_1} = 1 - \dfrac{T_2}{T_1}$
\end{itemize}

% ==================== 6. ЭЛЕКТРОСТАТИКА И ЭЛЕКТРОДИНАМИКА ====================
\section{Электростатика и Электрический ток}

\subsection{Электростатика}
\begin{itemize}
    \item \textbf{Закон Кулона:} $F = k \dfrac{|q_1 q_2|}{\epsilon r^2}$, \quad \textit{где} $k = \dfrac{1}{4\pi\epsilon_0} \approx 9 \cdot 10^9 \text{ Н}\cdot\text{м}^2/\text{Кл}^2$.
    \item \textbf{Напряженность точечного заряда:} $E = \dfrac{k |q|}{\epsilon r^2}$
    \item \textbf{Связь напряжения и напряженности в однородном поле:} $U = E d$
    \item \textbf{Потенциал точечного заряда:} $\varphi = \dfrac{k q}{\epsilon r}$
    \item \textbf{Емкость конденсатора:} $C = \dfrac{q}{U} = \dfrac{\epsilon \epsilon_0 S}{d}$
    \item \textbf{Энергия заряженного конденсатора:} $W_C = \dfrac{q U}{2} = \dfrac{C U^2}{2} = \dfrac{q^2}{2C}$
\end{itemize}

\subsection{Постоянный ток}
\begin{itemize}
    \item \textbf{Сила тока:} $I = \dfrac{q}{t}$
    \item \textbf{Сопротивление проводника:} $R = \rho \dfrac{l}{S}$
    \item \textbf{Закон Ома для участка цепи:} $I = \dfrac{U}{R}$
    \item \textbf{Закон Ома для полной цепи:} $I = \dfrac{\mathcal{E}}{R + r}$ \quad ($\mathcal{E}$ --- ЭДС, $r$ --- внутреннее сопротивление).
    \item \textbf{Закон Джоуля--Ленца (Выделяющееся тепло):} $Q = I^2 R t = U I t = \dfrac{U^2}{R} t$
\end{itemize}

% ==================== 7. МАГНЕТИЗМ И ЭЛЕКТРОМАГНИТНАЯ ИНДУКЦИЯ ====================
\section{Магнетизм и Электромагнитная индукция}

\begin{itemize}
    \item \textbf{Сила Ампера:} $F_A = B I l \sin\alpha$
    \item \textbf{Сила Лоренца:} $F_L = q v B \sin\alpha$
    \item \textbf{Радиус орбиты заряда в магнитном поле:} $R = \dfrac{m v}{q B}$
    \item \textbf{Магнитный поток:} $\Phi = B S \cos\alpha$
    \item \textbf{Закон электромагнитной индукции Фарадея:} $\mathcal{E}_{\text{инд}} = -\dfrac{\Delta\Phi}{\Delta t}$
    \item \textbf{ЭДС индукции в движущемся проводнике:} $\mathcal{E}_{\text{инд}} = v B l \sin\alpha$
    \item \textbf{Индуктивность и ЭДС самоиндукции:} $\Phi = L I, \quad \mathcal{E}_{\text{си}} = -L \dfrac{\Delta I}{\Delta t}$
    \item \textbf{Энергия магнитного поля катушки:} $W_M = \dfrac{L I^2}{2}$
\end{itemize}

% ==================== 8. КОЛЕБАНИЯ И ВОЛНЫ ====================
\section{Колебания и Волны}

\subsection{Механические и электромагнитные колебания}
\begin{itemize}
    \item \textbf{Уравнение гармонических колебаний:} $x(t) = A \cos(\omega t + \varphi_0)$
    \item \textbf{Математический маятник:} $T = 2\pi \sqrt{\dfrac{l}{g}}$
    \item \textbf{Пружинный маятник:} $T = 2\pi \sqrt{\dfrac{m}{k}}$
    \item \textbf{Колебательный контур (Формула Томсона):} $T = 2\pi \sqrt{L C}$
\end{itemize}

\subsection{Волны}
\begin{itemize}
    \item \textbf{Длина волны:} $\lambda = v T = \dfrac{v}{\nu}$
    \item \textbf{Скорость света в среде:} $v = \dfrac{c}{n}$ \quad ($n$ --- показатель преломления).
\end{itemize}

% ==================== 9. ОПТИКА ====================
\section{Оптика}

\begin{itemize}
    \item \textbf{Закон преломления света (Закон Снеллиуса):} $\dfrac{\sin\alpha}{\sin\beta} = \dfrac{n_2}{n_1}$
    \item \textbf{Предельный угол полного внутреннего отражения:} $\sin\alpha_{\text{пр}} = \dfrac{n_2}{n_1}$
    \item \textbf{Формула тонкой линзы:} $\pm \dfrac{1}{F} = \pm \dfrac{1}{d} \pm \dfrac{1}{f}$
    \item \textbf{Оптическая сила линзы:} $D = \dfrac{1}{F}$ (в диоптриях).
    \item \textbf{Дифракционная решетка:} $d \sin\varphi = m \lambda$
\end{itemize}

% ==================== 10. КВАНТОВАЯ И ЯДЕРНАЯ ФИЗИКА ====================
\section{Квантовая и Ядерная физика}

\begin{itemize}
    \item \textbf{Энергия и импульс фотона:} $E = h \nu = \dfrac{h c}{\lambda}, \quad p = \dfrac{h}{\lambda}$
    \item \textbf{Уравнение Эйнштейна для фотоэффекта:} $h \nu = A_{\text{выход}} + \dfrac{m v_{\max}^2}{2} = A_{\text{выход}} + e U_z$
    \item \textbf{Красная граница фотоэффекта:} $\nu_{\min} = \dfrac{A_{\text{выход}}}{h}, \quad \lambda_{\max} = \dfrac{h c}{A_{\text{выход}}}$
    \item \textbf{Дефект массы и энергия связи ядра:} $\Delta m = (Z m_p + N m_n) - M_{\text{ядра}}, \quad E_{\text{св}} = \Delta m \cdot c^2$
    \item \textbf{Закон радиоактивного распада:} $N(t) = N_0 \cdot 2^{-\frac{t}{T}}$ \quad ($T$ --- период полураспада).
\end{itemize}

% ==================== СВОДНАЯ ТАБЛИЦА ====================
\newpage
\section{Сводная таблица параметров вращательного движения}

\begin{table}[h!]
\centering
\caption{Взаимосвязь основных величин при движении по окружности}
\vspace{2mm}
\begin{tabular}{lcccc}
\toprule
\textbf{Величина} & \textbf{Через $T$} & \textbf{Через $\nu$} & \textbf{Через $\omega$} & \textbf{Через $v$ и $R$} \\
\midrule
Период ($T$) & $T$ & $\frac{1}{\nu}$ & $\frac{2\pi}{\omega}$ & $\frac{2\pi R}{v}$ \\
Частота ($\nu$) & $\frac{1}{T}$ & $\nu$ & $\frac{\omega}{2\pi}$ & $\frac{v}{2\pi R}$ \\
Угловая скорость ($\omega$) & $\frac{2\pi}{T}$ & $2\pi\nu$ & $\omega$ & $\frac{v}{R}$ \\
Центростремительное ускорение ($a_n$) & $\frac{4\pi^2 R}{T^2}$ & $4\pi^2\nu^2 R$ & $\omega^2 R$ & $\frac{v^2}{R}$ \\
\bottomrule
\end{tabular}
\end{table}

\end{document}
